\section{Conclusion}
\label{sec:conclusion}

This paper has addressed three questions about VRA compliance in algorithmic state legislative redistricting and answered each affirmatively.

First, the 42\% minority population threshold identified as a practical bright line for majority-minority district formation at congressional scale holds at state house scale. VRASection applied to state house chambers in the five covered states (Alabama, Georgia, Louisiana, Mississippi, South Carolina) produces majority-minority districts using the same threshold calibration as congressional-level analysis. The threshold holds because it is a geometric property of minority population aggregation that is scale-invariant: the break between achievable and unachievable majority-minority configurations occurs at approximately the same minority population share regardless of district size. South Carolina---the one covered state that fails the VRASection test at congressional scale---succeeds at state house scale because smaller districts can match the geographic extent of its dispersed Black population concentrations.

Second, the Callais disentanglement methodology applies at state legislative scale with identical mathematical structure and qualitatively identical results. The regression of district boundary decisions on minority population share versus partisan composition produces the same pattern at state house scale as at congressional scale: large, significant coefficients on minority population share and near-zero coefficients on partisan composition, confirming that the algorithm's boundary decisions in majority-minority districts are explained by minority community geography rather than partisan strategy. This categorical satisfaction of the Callais disentanglement requirement is a consequence of the algorithm's architectural inability to access partisan data.

Third, state house maps can achieve minority opportunity district targets more efficiently than congressional maps. The comparison across 13 states shows that state house VRASection produces 7.6 times as many majority-minority districts as congressional VRASection (378 versus 50) at substantially higher compactness levels (mean PP 0.365 versus 0.323 for majority-minority districts). The compactness penalty for majority-minority districts is smaller at state house scale (5.7\%) than at congressional scale (10.3\%), reflecting the more granular minority community matching possible with smaller district sizes.

The broader conclusion for algorithmic redistricting practice is that VRASection should be incorporated as a standard component of state legislative redistricting pipelines, not only congressional redistricting. States with substantial minority populations that may not trigger Section 2 obligations at congressional scale may nevertheless face those obligations at state house scale, because smaller districts lower the geographic compactness barrier to majority-minority district formation. The algorithm handles this expanded scope seamlessly: VRASection's threshold-based architecture scales from 4-district congressional maps to 400-district state house maps without modification, detecting minority community concentrations at the appropriate geographic scale for the chamber being drawn.

Future work should address three open questions. First, the empirical analysis here is based on 2020 census data; the results should be validated across the 2000 and 2010 census years, particularly for states where minority population geography has changed substantially over time. Second, the analysis treats the three \textit{Gingles} preconditions asymmetrically: we analyze precondition 1 (geographic compactness) empirically but take preconditions 2 (political cohesion) and 3 (majority bloc voting) as given. A full Section 2 analysis requires empirical analysis of actual voting data to establish cohesion and bloc voting, which is outside the scope of this paper. Third, the analysis focuses on Black and Hispanic majority-minority districts; Native American and Asian-Pacific Islander majority opportunity districts at state legislative scale are underexplored and warrant separate analysis.

The principal policy implication is encouraging: algorithmic redistricting at state legislative scale can satisfy Section 2 obligations as effectively as at congressional scale, with the additional advantage that the granular district scale allows more precise minority community boundary matching and a smaller compactness-VRA tradeoff. States adopting algorithmic redistricting for their congressional maps should consider simultaneous adoption for their state legislative maps, where the algorithm's performance advantages are, if anything, stronger.
